% !TEX root = thesis.tex
\chapter{Benchmarking Wireless Channel Representations via Downstream Tasks}\label{sec:tasks}

To evaluate the efficacy of the LWM and the Convolutional Autoencoder as universal feature extractors, we designed three distinct downstream applications. These tasks are selected to assess different capabilities of the learned representations: semantic understanding (classification), spatial reasoning (beam selection), and resource management (power allocation).

As previously discussed, in this work, we compare three representations:
\begin{enumerate}
    \item The \textbf{Latent Embeddings} produced by the LWM transformer. The LWM produces two different embeddings that will be compared separately:
    \begin{itemize}
        \item \textbf{Channel Embeddings}: A complete representation produced by embedding every patch of the channel.
        \item \textbf{CLS Embeddings}: A compact global representation of the channel.
    \end{itemize}
    \item The \textbf{Compressed Vector} $\mathbf{v}$ generated by the Autoencoder.
    \item The \textbf{Raw Channel State Information (CSI)}, used as the baseline.
\end{enumerate}

The representations are evaluated across three downstream tasks:
\begin{itemize}
    \item \textbf{Line-of-Sight (LoS) vs. Non-Line-of-Sight (NLoS) Classification}: Assessing the model's ability to capture physical propagation environments.
    \item \textbf{Beam Selection}: Evaluating spatial awareness and directionality in the latent space.
    \item \textbf{Power Allocation}: Testing the features' capability for multi-user resource optimization.
\end{itemize}

Additionally, the fine-tuning capabilities of the LWM are evaluated. For the CLS Embedding, we evaluate the performance \textit{without fine-tuning} (frozen encoder), \textit{with partial fine-tuning} (last three layers), and \textit{with full fine-tuning}.

Finally, we analyze the different methods using three metrics: \textit{Data Efficiency} (performance vs. training samples), \textit{Noise Robustness} (performance vs. SNR), and \textit{Computational Complexity} (FLOPs and inference time).

\section{Line-of-Sight (LoS) Classification}\label{sec:task_los_nlos}

The Line-of-Sight (LoS) classification is the simplest task, serving as a primary benchmark for feature quality. In wireless communications, distinguishing whether a user has a direct path to the BS is critical for beam management, handover decisions, and other tasks. 

This task evaluates whether the compressed representations preserve the physical characteristics of the propagation environment. If a model can distinguish LoS from NLoS solely from a compressed vector, it indicates that the encoder has successfully learned the underlying physics of the channel.

\subsection{Problem Formulation}\label{sec:task_los_nlos_problem}

The problem is modeled as a binary classification task. Given an input representation $\mathbf{x}$ (derived from the channel $\mathbf{H}$), the model predicts the propagation condition $y \in \{0, 1\}$, where:
\begin{itemize}
    \item $y = 1$ denotes a \textbf{Line-of-Sight (LoS)} condition.
    \item $y = 0$ denotes a \textbf{Non-Line-of-Sight (NLoS)} condition.
\end{itemize}

The dataset is generated using the DeepMIMO engine \cite{Alkhateeb2019}. The simulator utilizes a ray-tracing engine to extract CSI, user coordinates, and LoS labels from scenarios based on real-world city maps. To ensure numerical stability, all channel coefficients are scaled by a normalization factor $\alpha$.

\subsection{Downstream Model Architecture}\label{sec:task_los_nlos_model}

\begin{figure}[h]
    \centering
    \includegraphics[width=0.5\linewidth]{downstream_tasks/model_los_nlos_classification.pdf}
    \caption{Unified downstream MLP architecture for LoS/NLoS classification.}
    \label{fig:classification_head}
\end{figure}

For this task, the same MLP architecture is used for the three different representations. It consists of a linear transformation, followed by a non-linear activation and a layer that maps the hidden dimension into the number of classes, in this case 2.

\subsubsection{Input Adaptation}\label{sec:tasK_los_nlos_input}

Since the representations have different shapes, they are flattened before entering the MLP:
\begin{itemize}
    \item \textbf{LWM}: The input $\mathbf{E} \in \mathbb{R}^{(P+1) \times D}$ is flattened to a vector of size $(P+1)D$.
    \item \textbf{Autoencoder}: The latent vector $\mathbf{z} \in \mathbb{R}^{D_{\text{latent}}}$ is used directly.
    \item \textbf{Raw Channel}: The complex matrix $\mathbf{H} \in \mathbb{C}^{M \times N}$ is treated as two real channels and flattened to size $2MN$.
\end{itemize}

\subsection{Training}\label{sec:task_los_nlos_training}

The downstream model is trained to minimize the cross entropy

\begin{equation}
    L = - \frac{1}{B} \sum_{i=1}^{B} \left[ y_i \ln(\hat{y}_i) + (1 - y_i) \ln(1 - \hat{y}_i) \right]
\end{equation}

where $y_i$ is the ground-truth label, $\hat{y}_i$ is the predicted probability and $B$ is the batch size. 

\subsection{Evaluation Metrics}\label{sec:task_los_nlos_evaluation}

Since the dataset may contain imbalances depending on the simulated environment (e.g., more NLoS users in urban canyons), we employ the \textbf{Weighted F1-Score} as the primary metric, alongside standard Accuracy.

\begin{equation}
    F1_{\text{weighted}} = \frac{N_1}{N} F1_1 + \frac{N_0}{N} F1_0
\end{equation}


$N$ represents the total number of samples and $N_i$ is the true number of instances in each class. The Weighted F1-Score calculates the metric for each class independently and computes the average weighted by the number of true instances for each class, providing a robust measure of performance. 

\begin{equation}
    F1_i = 2 \cdot \frac{\text{Precision}_i \cdot \text{Recall}_i}{\text{Precision}_i + \text{Recall}_i}
\end{equation}

\section{Beam Selection}\label{sec:task_beam}

The second downstream task evaluates a critical problem in Massive MIMO systems: beam selection. In high frequency systems with multiple users, beamforming becomes fundamental to overcome the path loss problem. 

This task evaluates the model's capability to understand \textbf{spatial relationships}. We want to understand whether the compressed representation $\mathbf{v}$ and the embeddings captures enough geometrical information to be used as input to perform beam selection.

\subsection{Problem Formulation}\label{sec:task_beam_problem}

The problem is modeled as a multi-class classification task. In this simulation the BS is equipped with a Uniform Linear Array (ULA) serving a single-antenna user. The BS uses a pre-defined codebook

\begin{enumerate}
    \item \textbf{Angles}: The codebook creates $K_{\text{beams}}$ angles spaced evenly between a \textbf{Field of View (FoV)}.
    \begin{equation}
            \phi_k = - \frac{\text{FoV}}{2} + k \cdot \Delta \phi , \quad k \in {0, ..., K-1}
    \end{equation}
    where $\phi_k$ is the $k$-th angle and $\Delta \phi = \frac{FoV}{K-1}$
    \item \textbf{Steering vector}: Consider a ULA with $M$ antennas and inter-element spacing $d$. We define the steering vector $\mathbf{w} \left( \phi_k\right)$ as
    \begin{equation}
            \mathbf{w} \left( \phi_k\right) = \frac{1}{\sqrt{M}} \left[ 1, e^{-j \frac{2 \pi}{\lambda} d \sin \left( \phi_k \right) }, e^{-j \frac{2 \pi}{\lambda} d (2) \sin \left( \phi_k \right) }, ..., e^{-j \frac{2 \pi}{\lambda} d \left( M-1 \right) \sin \left( \phi_k \right) }\right] ^{\top}
    \end{equation}
    \item \textbf{Normalization}: $\mathbf{w} \left( \phi_k\right)$ is divided by $|| \mathbf{w} \left( \phi_k\right) ||_{2}$ to ensure constant transmission to every beam index.
    \item \textbf{Codebook Matrix}: The steering vectors are stacked to create the Codebook Matrix $\mathbf{F} \in \mathbb{C}^{K \times M}$
    \begin{equation}
            \mathbf{F} = \begin{bmatrix} \mathbf{w}(\phi_0)^H \\ \mathbf{w}(\phi_1)^H \\ \vdots \\ \mathbf{w}(\phi_{K-1})^H \end{bmatrix}
    \end{equation}
    \item \textbf{Subcarrier averaging}: Since $\mathbf{H}$ has shape ($M \times N$), we average the subcarriers to create an effective spatial channel vector $\bar{\mathbf{h}}_i \in \mathbb{C}^{M \times 1}$
    \begin{equation}
            \bar{\mathbf{h}}_i = \frac{1}{N} \sum_{n=1}^{N} \mathbf{h}_{i,n}
    \end{equation}
    \item \textbf{Beam Selection}: We select the best beam by selecting the $k^{*}$ that maximizes the power received by the user. This value is the label for the selected user.
    \begin{equation}
            k^{*} = \operatorname*{argmax}_{k \in \left\{ 0, \cdot, K-1\right\}} \left| \mathbf{w}(\phi_k)^H  \bar{\mathbf{h}}_i \right|^2
    \end{equation}
\end{enumerate}

\subsection{Downstream Architecture}\label{sec:task_beam_model}

To ensure a fair comparison, a unified downstream architecture is employed for all representations. We utilize a standardized MLP acting as a classification head.

The network consists of a sequence of three fully connected blocks. Each block performs the following operations:
\begin{equation}
    \mathbf{h}_{l+1} = \text{Dropout}\left( \sigma \left( \text{BN} \left( \mathbf{W}_l \mathbf{h}_l + \mathbf{b}_l \right) \right) \right)
\end{equation}
where $\text{BN}$ denotes Batch Normalization, $\sigma$ is the ReLU activation function, and the dropout rate is set to $D_{\text{rate}}$. The dimensionality of the features is halved at each successive block to force the model to distill discriminative features. The final layer maps the features to the $K$ class logits.

\begin{figure}[h]
    \centering
    \includegraphics[width=0.75\linewidth]{downstream_tasks/model_beam_selection.pdf}
    \caption{Unified downstream MLP architecture for Beam Selection.}
    \label{fig:custom_classification_head}
\end{figure}

\subsection{Training}\label{task_beam_training}

In this task, the model is trained to minimize the cross entropy for $K$ classes.

\begin{equation}
    L = - \sum_{k=1}^{K} y_k \ln(\hat{y}_k)
\end{equation}

It is also necessary to perform the same input adaptations as in the previous task. Finally, the task is evaluated by the same metric, Weighted F1-Score.

\section{Power Allocation}\label{sec:task_power}

The final downstream task evaluates the performance of the learned representation in a \textbf{resource allocation} task. Specifically, we focus on a Multi-User MIMO (MU-MIMO) system. Unlike the previous tasks, this one is a regression problem where the goal is to maximize the system's spectral efficiency under total power budget constraint.

The task verifies if the representation is enough to deal with interference patterns between users.

\subsection{Problem Formulation}\label{sec:task_power_problem}

Consider a system with $K$ single-antenna users served at the same time by a BS with $M$ antennas. The received signal by the user $k$ is

\begin{equation}
    y_k = \mathbf{h}_k^H \mathbf{w}_k \sqrt{p_k} s_k + \sum_{j \neq k} \mathbf{h}_k^H \mathbf{w}_j \sqrt{p_j} s_j + n_k
\end{equation}

where:
\begin{itemize}
    \item $\mathbf{h}_k \in \mathbb{C}^{M \times 1}$ is the channel vector of the $k$-th user
    \item $\mathbf{w}_k \in \mathbb{C}^{M \times 1}$ is the beamforming vector of the $k$-th user
    \item $s_k$ is the symbol transmitted to user $k$ with $\mathbb{E} [|s_k|^2] = 1$.
    \item $p_k$ is the allocated power for the $k$-th user
    \item $\mathbf{h}_k^H \mathbf{w}_k \sqrt{p_k} s_k$ represents the signal for the $k$-th user
    \item $\sum_{j \neq k} \mathbf{h}_k^H \mathbf{w}_j \sqrt{p_j} s_j $ represents the interference suffered by the $k$-th user.
    \item $n_k$ is Additive white Gaussian noise (AWGN) $\mathcal{CN}(0, \sigma^2)$
\end{itemize}

For the beam directions, we assume Maximum Ratio Transmission (MRT) precoding \cite{mrt1999}. The main reason to choose this strategy is due to its low computational complexity compared to strategies that actively cancel interference, like Zero-forcing (ZF) \cite{spencer2004zf}. Also, the most complicated part of the interference management is done by the power allocation.

\begin{equation}
    \mathbf{w}_k = \frac{\mathbf{h}_k}{\| \mathbf{h}_k \|}
\end{equation}

\subsection{Downstream Model Architecture}\label{sec:task_power_model}

Similar to the previous tasks, we use a unified downstream architecture to ensure a fair comparison. A 1D Residual Convolutional Neural Network (ResCNN) was employed in the power allocation. The main reason to choose a more complex model is that we were unable to achieve consistent results using a simpler one, since this task requires managing the inter-user interference.

The forward pass begins with a one-dimension (1D) convolution entry block and a non-linear activation. The core of the model is a sequence of residual blocks. Each one applies two layers of one dimension convolution with batch normalization and non-linear activation. We add a residual connection to mitigate the vanishing gradient problem \cite{he2016deep}.

In the final block, the model matches one power logit per user. To ensure the $P_{\text{total}}$ budget constraint, a Softmax activation is applied.

\begin{figure}[ht]
    \centering
    \includegraphics[width=0.5\linewidth]{downstream_tasks/model_power_allocation.pdf}
    \caption{Unified downstream MLP architecture for power allocation.}
    \label{fig:regression_head}
\end{figure}

In the last block, the output is mapped into $K$ users and the last activation function is a softmax to enforce the power budget constraint equal to $P_{\text{total}}$

The input adaptations are the same used in the previous tasks, but accounting for the $K$.

\begin{itemize}
    \item \textbf{LWM}: Input flattened into $K$ vectors of size $(P + 1) D$.
    \item \textbf{Autoencoder}: Latent vector $z$ is transformed into a sequence of $K$ vector with size $D_{\text{latent}}$
    \item \textbf{Raw Channel}: The complex matrix $\mathbb{H}$ is treated as two real channels and flattened to $K$ vectors with size $2 M N $.
\end{itemize}

\subsection{Training}\label{sec:task_power_training}

Optimizing the spectral efficiency using Deep Learning is a difficult task due to the non-convex nature of the function. To address this problem, we choose to implement a hybrid training strategy with a fixed SNR equal to 5 dB during the training.

\subsubsection{Supervised Warm-up}\label{sec:task_power_supervised}

In the first phase of the training ($60\%$ of the epochs), the model is trained using a supervised strategy to learn with a classical optimization algorithm. We generate power allocation $\mathbf{p}^{*}$ labels for a subset of training dataset using \textbf{Projected Gradient Descent (PGD)} \cite{boyd2004convex}.

Specifically, the PGD algorithm iteratively maximizes the sum-rate capacity of the multi-user channel. At each iteration, it computes a gradient ascent step to increase the system throughput. To ensure the physical constraints, the algorithm checks if all power labels are non-negative ($p_k \geq 0$) and the total power of the system is less than the power budget ($\sum_{k=1}^{K} p_k \leq P_{\text{total}}$).

The model is trained to minimize the Mean Squared Error (MSE) between the predicted powers $\hat{\mathbf{p}}$ and the target.

\begin{equation}
    \mathcal{L}_{\text{warmup}} = \|\hat{\mathbf{p}} - \mathbf{p}^*\|^2
\end{equation}

This phase is very important because, due to the non-convex nature of the problem and the random initialization of the weights, the model can be stuck in a local minimum. The warm-up phase stabilizes training and provides a good initialization.

\subsubsection{Unsupervised learning}\label{sec:task_power_unsupervised}

After the warm-up, the model begins an unsupervised learning phase with the objectives of maximizing the spectral efficiency. Since the standard optimizers want to minimize loss, the negative SE is defined as loss:

\begin{equation}
    \text{SE}(\hat{\mathbf{p}}) = - \frac{1}{B} \sum_{k=1}^{K} \log_2 \left( 1 + \text{SINR}_k(\hat{\mathbf{p}}) \right)
\end{equation}

where $B$ is the number of samples in the batch and $\text{SINR}_k(\hat{\mathbf{p}})$ is the Signal-to-Interference-plus-Noise Ratio for the $k$-th user for a given set of power $\hat{\mathbf{p}}$ for each of the $k$ users.

\subsection{Dataset Generation: User Grouping}\label{sec:task_power_dataset}

An important aspect of MU-MIMO systems is the user selection. Randomly sampling $K$ users could result in groups with highly correlated channels, making it impossible to manage the interference through power allocation, or very different path losses, which could result in a dominant user receiving most of the power.
To ensure a meaningful benchmark, we implemented a greedy sampling algorithm based on two criteria:

\begin{enumerate}
    \item \textbf{Spatial Correlation}: Users must have spatial correlation $\rho_{\text{min}} \leq \rho \leq \rho_{\text{max}}$.
    \item \textbf{Gain Ratio}: To minimize the near-far problem, the ratio between the weakest and strongest user should be less than $\gamma_{max}$.
\end{enumerate}

\subsection{Evaluation Metrics}\label{sec:task_power_evaluation}

To evaluate the performance of the downstream model based on the learned representation, we compare the spectral efficiency of predicted powers against three baselines:

\begin{enumerate}
    \item \textbf{Equal Power Allocation (EPA)}: The lower boundary of our system, where the power is shared equally among users $p_k = P_{\text{total}} / K $.
    \item \textbf{Iterative PGD (Optimal)}: A numerical solver that runs for $i$ iterations with $r$ different random initializations.
    \begin{equation}
        \begin{aligned}
            \max_{\mathbf{p}} \quad & \sum_{k=1}^{K} \log_2(1 + \text{SINR}_k) \\
            \text{s.t.} \quad & \sum_{k=1}^{K} p_k \leq P_{\text{total}} \\
            & p_k \geq 0, \quad \forall k \in \{1, \dots, K\}
        \end{aligned}
        \label{eq:sum_rate_optimization}
    \end{equation}
    
    where $\mathbf{p} = [p_1, p_2, \dots, p_K]^T$ is the allocated power vector and

    \begin{equation}
        \text{SINR}_k = \frac{|\mathbf{h}_k^H \mathbf{w}_k|^2 p_k}{\sum_{j \ne k} |\mathbf{h}_k^H \mathbf{w}_j|^2 p_j + \sigma^2}
    \end{equation}
\end{enumerate}